\documentstyle[% 


righttag% 


,11pt% 


,amssymb% 


,russian%               remove in Eng. 


,izvru%                 Set izven% in Eng. 


]{amsart} 


 


%       Math definitions 


 


\newcommand{\tts}[1]{} 


\renewcommand{\baselinestretch}{1.03}


 


%\hoffset=-15mm%                  For Eng. 


  


\setcounter{page}{43} 


\god{1999} 


\nomer{10~(449)} %                        Remove in Eng. 


%\nomer{Vol.\,43, No.\,10}%               Set in Eng. 


%\str{\pageref{begin}--\pageref{end}}%  Set in Eng. 


\udk{517.982} 


 


\author{\it Г.С.\,Тюгашов} 


% NB:   ^^ remove in Eng. 


\title{$k$-растянутость топологических пространств\\ и выпуклость} 


 


\date{} 


%\pagestyle{myheadings}%         Set in Eng. 


\pagestyle{plain} %              Remove in Eng. 


\begin{document} 


%\thispagestyle{plain}%          Set in Eng. 


\maketitle 


\label{begin} 


 


Приведенные ниже утверждения обобщают результаты, 


полученные в [1] и [2] для 


$\alpha$-растянутых топологических пространств.  


 


Напомним некоторые обозначения и определения из 


[3], где (как и в данной статье) все топологические  


пространства предполагаются $T_1$-пространствами. 


 


Если $<$ --- линейное упорядочение на пространстве $X$, то (по 


определению) $X_{<x}\,(X_{\le x})\,=\{y\in X:y<x\ (y\le x)\}$. Если  


$<$ --- линейное упорядочение на топологическом пространстве $X$, $P$ 


--- класс топологических пространств и для любого $x\in X$ 


имеем $[X_{<x}]\setminus X_{<x}\in P$, то линейное упорядочение $<$ 


называется  


слабо $P$-согласованным с топологией пространства $X$ или 


слабо $P$-согласованным на $X$. Если, кроме того, для любого $x\in X$ 


при $[X_{<x}]\setminus X_{<x}\ne\varnothing$ имеем 


$x\in[X_{<x}]\setminus X_{<x}$, то линейное  


упорядочение $<$ назовем $P$-согласованным с топологией  


пространства $X$ ($P$-согласованным на $X$). 


\medskip 


 


{\bf Определение.} Пусть $P$ --- произвольный класс  


топологических пространств. Топологическое пространство $X$  


назовем (слабо) $P$-растянутым ($P$-левым), если на $X$ существует  


(слабо) $P$-согласованное линейное упорядочение (вполне  


упорядочение); если на $X$ существует (слабо) $P$-согласованное 


упорядочение, обратное к вполне упорядочению, то пространство $X$  


назовем (слабо) $P$-правым. 


\medskip 


 


Если класс $P$ состоит из всех топологических пространств,  


включающих не более $k$ точек, где $k$ --- любое натуральное число, то  


основное определение становится определением (слабо) 


$k$-растянутых, (слабо) $k$-левых и (слабо) $k$-правых пространств; 


если при этом $k=1$, то получаем определение (слабо) 


$\alpha$-растянутых, (слабо) $\alpha$-левых, (слабо) $\alpha$-правых 


пространств. Наконец, если 


класс $P$ состоит из всех пустых топологических пространств, то  


основное определение становится определением растянутых, левых,  


правых пространств. 


 


Пусть далее $<$ --- линейное упорядочение на топологическом  


пространстве $X$ и $A\subset X$. Говорят, что множество $A$ 


выпукло, если  


из $x\in A$, $y\in A$ и $x\le z\le y$ следует, что $z\in A$.  


 


\begin{thm} Если на топологическом пространстве $X$ 


существует $k$-согласованное линейное упорядочение $<$,  


удовлетворяющее у\,с\,л\,о\,в\,и\,ю $Z:$


$$


\text{для $\forall\,x\in X$ множество 


$[X_{<x}]\setminus X_{<x}$ выпукло},


$$


то любое бикомпактное подмножество пространства $X$  


имеет минимальный элемент относительно упорядочения $<$. 


\end{thm} 


 


{\bf Доказательство.} Для любого $x\in A$ положим $A_x=A\cap[X_{<x}]$. 


Так как линейное упорядочение $<$\; $k$-согласовано на $X$, то 


$x\in A_x$ 


для любого $x\in A$ и семейство $\{A_x:x\in A\}$ является центрированной  


системой замкнутых множеств в бикомпакте $A$ и потому имеет  


непустое пересечение $B=\cap\{A_x:x\in A\}$. Покажем, что $|B|\le k$. 


Предположим (от противного), что $|B|>k$, и зафиксируем $k+1$ 


точек $x_1,\dotsc,x_k,x_{k+1}\in B$,


$x'=\min\{x_i:i=1,\dotsc,k,k+1\}$. Тогда  


$x_i\in[X_{<x'}]\setminus X_{<x'}$ для любого $i=1,\dotsc,k,k+1$, 


следовательно, 


$|[X_{<x'}]\setminus X_{<x'}|>k$, что противоречит $k$-согласованности  


упорядочения $<$ на $X$. Значит, $|B|\le k$ и во множестве $B$ 


существует  


минимальный элемент $x_0$. 


 


Предположим, что в $A$ нет минимального элемента и выберем  


$x_k<\dotsb<x_1<x_0\in A$. Так как $<$ --- $k$-согласованное линейное  


упорядочение на $X$, то $x_k\in[X_{<x_k}]\setminus X_{<x_k}$, 


а из $x_0\in B$ следует, что 


$x_0\in[X_{<x_k}]\setminus X_{<x_k}$, поэтому в силу условия $Z$ 


получаем, что 


$x_R\in[X_{<x_k}]\setminus X_{<x_k}$ для любого $R=0,1,\dotsc,k$. Это 


противоречие 


доказывает теорему. 


 


\begin{exam} Пусть $X=\{1/n:n\in N\}\cup\{0\}\subset R$. Топология на  


$X$ индуцирована естественной топологией вещественной прямой $R$.  


Для любого $x\in X$, $x\ne0$, положим $x<0$, а для остальных пар 


$x,y\in X$, 


$x\ne0$, $y\ne0$, порядок $<$ совпадает с естественным порядком 


вещественной прямой. Тогда 


$|[X_{<x}]\setminus X_{<x}|=1$ для любой  


точки $x\in X$, следовательно, выполнено условие $Z$ и $<$ --- слабо 


$\alpha$-согласованное линейное упорядочение на $X$ (но не 


$\alpha$-согласованное). В бикомпакте $X$ нет минимального элемента  


относительно упорядочения $<$. 


\end{exam} 


 


\begin{exam} Пусть $X=[0,1]\subset R$. Топология на $X$ индуцирована  


естественной топологией вещественной прямой $R$. Для любого 


$x\in X$, $x\ne0$, положим $x<0$, а для остальных пар $x,y\in X$, 


$x\ne0$, $y\ne0$, 


линейный порядок $<$ совпадает с естественным порядком  


вещественной прямой. Тогда 


$|[X_{<x}]\setminus X_{<x}|\le 2$ и 


$x\in[X_{<x}]\setminus X_{<x}$ для 


любой точки $x\in X$, следовательно, $<$ --- 2-согласованное линейное  


упорядочение на $X$. Очевидно, условие $Z$ не выполнено. В  


бикомпакте $X$ нет минимального элемента относительно  


упорядочения $<$. 


\end{exam} 


 


\begin{exam} Пусть $X$ --- множество рациональных чисел с 


естественной топологией и $<$ --- обратный порядок к минимальному  


вполне упорядочению на $X$. Так как множество $X$ счетно, то для  


любого $x\in X$ множество $[X_{<x}]\setminus X_{<x}$ конечно и 


$[X_{<x}]=X$, 


следовательно, выполнено условие $Z$ и $<$ является 


$P$-согласованным линейным упорядочением на $X$, где $P$ --- класс всех  


конечных топологических пространств. Очевидно, упорядочение $<$  


не является $k$-согласованным на $X$ ни для какого натурального $k$.  


Если в качестве подпространства $A$ в $X$ рассмотреть  


бикомпакт из примера 1, то в $A$ нет минимального элемента  


относительно упорядочения $<$. 


\end{exam} 


 


Примеры 1--3 показывают, что все посылки теоремы 1 являются  


существенными. 


 


\begin{thm} 


Для любого топологического пространства $X$  


следующие утверждения рав\-но\-сильны$:$ 


\begin{itemize} 


\item[1.] пространство $X$ слабо $k$-растянуто$;$ 


\item[2.] пространство $X$\; $k$-растянуто$;$ 


\item[3.] пространство $X$\; $k$-растянуто и удовлетворяет условию $Z$.  


\end{itemize} 


\end{thm} 


 


{\bf Доказательство.} Очевидно, $3)\Rightarrow 2)$, $2)\Rightarrow 1)$. 


Покажем, что 


$1)\Rightarrow 3)$. Пусть $<$ --- слабо $k$-согласованное линейное 


упорядочение  


на $X$. 


 


$1^\circ$. Положим $X'=\{x\in X:|[X_{<x}]\setminus X_{<x}|\ge k\}$. 


Без ограничения  


общности считаем $X\ne\varnothing$ (в противном случае полагаем  


$k=\max\{|[X_{<x}]\setminus X_{<x}|:x\in X\}$). 


 


Рассмотрим разбиение $F$ множества $X'$ на непересекающиеся  


подмножества, полагая, что точки $x,y\in X'$ принадлежат одному и  


тому же подмножеству $A\in F$, если 


$\min\{[X_{<x}]\setminus X_{<x}\}= 


\min\{[X_{<y}]\setminus X_{<y}\}$ 


($\min$ берется относительно упорядочения $<$). Пусть  


$x_A=\{[X_{<x}]\setminus X_{<x}\}$ для некоторого $x\in A\in F$, тогда 


$x_A\in A$ для любого  


$A\in F$. Очевидно, что каждый элемент семейства $F$ является  


выпуклым подмножеством множества $X$. Определим на множестве  


$X$ новый линейный порядок $\ll$ следующим образом: для любого  


$A\in F$, для любого $x\in A$, $x\ne x_A$, полагаем $x_A\ll x$; для 


всех  


остальных пар точек $x,y\in X$ порядок $\ll$ совпадает с порядком $<$.  


Тогда для любого $x\in X\setminus X'$ имеем $X_{\ll x}=X_{<x}$ и, 


следовательно,  


$|[X_{\ll x}]\setminus X_{\ll x}|<k$. Кроме того, если $x\in A\in F$, 


$x\ne x_A$, то $|[X_{\ll x}]\setminus X_{\ll x}|< 


|[X_{<x}]\setminus X_{<x}|=k$. 


 


Наконец, если $x=x_A$ для некоторого $A\in F$, то либо $x\notin[X_{\ll 


x}]$ и,  


значит, $|[X_{\ll x}]\setminus X_{\ll x}|<k$, либо $x\in[X_{\ll x}]$ и 


только в этом случае  


может быть $|[X_{\ll x}]\setminus X_{\ll x}|\ge k$. Переобозначим 


линейное  


упорядочение $\ll$ через $<$. Тогда на следующем шаге имеем слабо 


$k$-согласованное линейное упорядочение $<$ на пространстве $X$ 


такое, что, если $|[X_{<x}]\setminus X_{<x}|\ge k-1$ для некоторой точки 


$x\in X$, то  


$[X_{<x}]\setminus X_{<x}=B\cup C$, где $B$ --- выпуклое относительно 


$<$ подмножество в $X$, $|B|\le 1$, $x\in B$ при $B\ne\varnothing$, 


$|C|=k-1$. 


 


$2^\circ$. $X'=\{x\in X:|[X_{<x}]\setminus X_{<x}|\ge k-1\}$. 


Как и выше, рассмотрим  


разбиение $F$ множества $X'$ на непересекающиеся подмножества,  


полагая, что точки $x,y\in X'$ принадлежат одному и тому же  


подмножеству $A\in F$, если $\min\{C_1\}=\min\{C_2\}$ в представлениях 


$[X_{<x}]\setminus X_{<x}=B_1\cup C_1$, $[X_{<y}]\setminus 


X_{<y}=B_2\cup C_2$. 


 


Пусть $x_A=\min\{C\}$ и $[X_{<x}]\setminus X_{<x}=B\cup C$ 


для некоторого  


$x\in A\in F$. Тогда $x_A\in A$ для любого $A\in F$. Каждый элемент 


семейства  


$F$ является выпуклым подмножеством множества $X$. Определим на  


множестве $X$ новый линейный порядок $\ll$ следующим образом: 


для любого $A\in F$, для любого $x\in A$, $x\ne x_A$, полагаем $x_A\ll 


x$; для  


всех остальных пар точек $x,y\in X$ порядок $\ll$ совпадает с  


порядком $<$. Тогда для любого $x\in X\setminus X'$ имеем $X_{\ll 


x}=X_{<x}$ и,  


следовательно, $|[X_{\ll x}]\setminus X_{\ll x}|<k-1$. Кроме того, если 


$x\in A\in F$, 


$x\ne x_A$, то $|[X_{\ll x}]\setminus X_{\ll x}|<|[X_{<x}]\setminus 


X_{<x}|$ и $[X_{\ll x}]\setminus X_{\ll x}=B\cup C$, где 


$B$ --- выпуклое относительно $\ll$ подмножество в $X$, $|B|\le 1$, 


$x\in B$ 


при $B\ne\varnothing$, $|C|=k-2$. Наконец, если $x=x_A$ для некоторого 


$A\in F$, то  


либо $x\notin[X_{\ll x}]$ и, значит, 


$[X_{\ll x}]\setminus X_{\ll x}=B\cup C$, где $B=\varnothing$ 


(следовательно, выпукло), $|C|\le k-2$, либо $x\in[X_{\ll x}]$ и 


$[X_{\ll x}]\setminus X_{\ll x}=B\cup C$, где $|B|\le 2$, $B$ --- 


выпуклое относительно $\ll$ 


подмножество в $X$, $|B|\le 2$, $|C|\le k-2$. 


 


Итак, получено слабо $k$-согласованное линейное  


упорядочение $\ll$ на пространстве $X$, для которого при условии  


$|[X_{\ll x}]\setminus X_{\ll x}|\ge k-2$ (для некоторой точки $x\in X$)


справедливо  


представление $[X_{\ll x}]\setminus X_{\ll x}=B\cup C$, где $B$ --- 


выпуклое относительно  


$\ll$ подмножество в $X$, $|B|\le 2$, $x\in B$ при $B\ne\varnothing$, 


$|C|=k-2$. 


 


Продолжая описанную процедуру, на последнем $k$-м шаге  


получим слабо $k$-согласованное линейное упорядочение $\ll$ на  


пространстве $X$ такое, что если $|[X_{\ll x}]\setminus X_{\ll x}|\ge 0$ 


для некоторой  


точки $x\in X$, то $[X_{\ll x}]\setminus X_{\ll x}=B\cup C$, где $B$ 


--- выпуклое  


относительно $\ll$ подмножество в $X$, $|B|\le k$, $x\in B$ при 


$B\ne\varnothing$, $|C|=0$. 


 


Так как всякое подпространство слабо $k$-растянутого 


топологического пространства $k$-растянуто, то в силу теоремы 2  


справедливо 


 


\begin{cor} 


Любое подпространство $k$-растянутого 


топологического пространства $k$-рас\-тя\-нуто. 


\end{cor} 


 


\begin{thm} 


Если $f:X\to Y$ --- замкнутое бикомпактное  


$($необязательно непрерывное$)$ отображение топологического 


пространства $X$ на топологическое пространство $Y$ и пространство $X$ 


является 


$k$-растянутым, то $k$-растянуто и пространство $Y$. 


\end{thm} 


 


{\bf Доказательство.} Пусть $<$ есть $k$-согласованное линейное 


упорядочение на пространстве $X$ с условием $Z$. Для любого $y\in Y$ 


положим $x_y=\min f^{-1}(y)$ (такой минимум существует по теореме 1 в  


силу бикомпактности отображения $f$). Определим линейное 


упорядочение $\ll$ на пространстве $Y$ следующим образом: $y'\ll y''$ 


для $y'\ne y''\in Y$, если и только если $x_{y'}<x_{y''}$. В силу 


теоремы 2  


достаточно показать, что линейное упорядочение $\ll$ слабо  


$k$-согласовано на $Y$. 


 


Пусть $y^*\in Y$ и $x^*=x_{y^*}$, тогда $f(X_{<x^*})=Y_{\ll y^*}$. 


Действительно,  


если $y\ll y^*$, то $x_y<x^*$, следовательно, $x_y\in X_{<x^*}$ и 


$y=f(x_y)\in f(X_{<x^*})$. 


Значит, $Y_{\ll y^*}\subset f(X_{<x^*})$. Если $x<x^*$, то $x_{f(x)}\le 


x<x^*=x_{y^*}$ и, значит,  


$f(x)\ll y^*$, т.е. $f(x)\in Y\ll y^*$, следовательно, 


$f(X_{<x^*})\subset Y_{\ll y^*}$. Итак, 


$f(X_{<x^*})=Y_{\ll y^*}$. Тогда в силу замкнутости отображения $f$ 


имеем  


$[Y_{\ll y}]\setminus Y_{\ll y^*}\subset f([X_{<x^*}])\setminus Y_{\ll 


y^*}=f([X_{<x^*}])\setminus f(X_{<x^*})\subset f([X_{<x^*}])\setminus 


X_{<x^*}$ 


и, значит, $|[Y_{\ll y^*}]\setminus Y_{\ll y^*}|\le|[X_{<x^*}]\setminus 


X_{<x^*}|\le k$, что и требовалось 


доказать. 


 


Аналогично обоснованию теоремы 3 доказывается  


 


\begin{thm} 


Если $f:X\to Y$ --- открытое бикомпактное 


отображение на топологическое пространство $Y$ и на пространстве 


$X$ существует слабо $k$-согласованное линейное упорядочение такое, 


что обратный порядок $l$-согласован и удовлетворяет условию $Z$, то 


пространство $Y$ является $k$-растянутым. 


\end{thm} 


 


\begin{thm} 


Если $X$ -- не пустой $k$-растянутый бикомпакт  


счетной тесноты, то в $X$ существует точка счетного характера. 


\end{thm} 


 


{\bf Доказательство.} 


Пусть $x(x,X)\ge\aleph$ для любого $x\in X$ и $<$ --- 


$k$-согласованное линейное упорядочение на $X$, удовлетворяющее  


условию $Z$, $F$ --- семейство всех непустых бикомпактов счетного  


характера в $X$. Определим $F_\alpha\in F$ и $x_\alpha\in X$ для любого 


$\alpha<\aleph$ по 


правилу: $F_0=X$ и $x_0=\min X$. Пусть $\beta<\aleph_1$ и $F_\alpha\in 


F$, $x_\alpha\in X$ 


определены для всех $\alpha<\beta$ с условиями: a)~если 


$a'<\alpha''<\beta$, то  


$F_{\alpha''}\subset 


F_{\alpha'}\setminus[\{x_\alpha:\alpha\le\alpha'\}]$; б)~$x_\alpha=\min 


F_\alpha$ для любого $\alpha<\beta$. 


Тогда $\Phi_\beta=\cap\{F_\alpha:\alpha<\beta\}\in F$. Для 


$y_\beta=\min\Phi_\beta$ имеем  


$x_\alpha\le y_\beta$ любого $\alpha<\beta$, поэтому 


$[\{x_\alpha:\alpha<\beta\}]\subset[\{x\in X:x<y_\beta\}]$. 


Значит, $|[\{x_\alpha:\alpha<\beta\}]\subset\Phi_B|\le K$. Так как 


$\aleph(x,X)\ge\aleph_1$ для  


любого $x\in X$, то $|\Phi_\beta|<K$. Поэтому существует $F_\beta\in F$ 


такое, что  


$F_\beta\subset\Phi_\beta\setminus[\{x_\alpha:\alpha<\beta\}]$ и 


$x_\beta=\min F_\beta$. 


Тогда $\{x_\alpha:\alpha<\aleph_1\}$ --- свободная последовательность в 


$X$  


длины $\aleph_1$. Действительно, для любого $\beta<\aleph_1$ имеем  


$[\{x_\alpha:\alpha<\beta\}]\cap F_\beta=\varnothing$ и 


$[\{x_\alpha:\beta\le\alpha\}]\subset[F_\beta]$. 


Но в $X$ не может  


быть свободной последовательности длины $\aleph_1$, что и доказывает  


теорему. 


 


\begin{thm} 


Если $X$ есть $k$-растянутое, счетно компактное,  


регулярное $T_1$-пространство, то существуют такая точка $x\in X$ и  


такая последовательность $\eta=\{U_n:n\in N^+\}$ непустых открытых в 


$X$ 


множеств, что $\eta$ сходится к $x$. 


\end{thm} 


 


Для доказательства теоремы 6 сформулируем и докажем две  


леммы. Обозначим через $P$ семейство всех непустых открытых в $X$  


множеств и пусть $<$ --- $k$-согласованное линейное упорядочение на  


$X$, удовлетворяющее условию $Z$. 


 


\begin{lem} 


Если $U\in P$ и $|U|>1$, то существуют такие $V\in P$ и  


$x\in U$, что $V\subset U$ и $y<x$ для любого $y\in V$. 


\end{lem} 


 


{\bf Доказательство} (от противного). Пусть существует $U\in P$, 


$|U|>1$, для которого не выполнено заключение леммы 1, тогда в $U$  


нет наибольшего элемента, иначе $V=U\setminus\{\max U\}$ 


искомое. Положим  


$P^*=\{V\in P:V\subset U\}$. Из сделанного предположения следует, что  


для любого $x\in U$ и каждого $V\in P^*$ получим 


$V\setminus[X_{<x}]=\varnothing$. 


Определим последовательности $\eta^i=\{V^i_j:j\in N^+\}\subset P^*$ и 


$\xi^i=\{x^i_j:j\in N^+\}\subset U$, 


где $i\in\{1,2,\dotsc,k,k+1\}$. 


 


Выберем $V^i_1\in P^*$ так, чтобы было: 


a)~$V^i_1\cap V^j_1=\varnothing$ при 


$i,j\in\{1,2,\dotsc,k,k+1\}$, $i\ne j$ и б)~$V^i_1\subset U$, 


$i\in\{1,2,\dotsc,k,k+1\}$. 


 


Фиксируем любые $x^i_1\in V^i_1$ для каждого 


$i\in\{1,2,\dotsc,k,k+1\}$. 


Пусть $l\in N^+$ и $V^i_l\in U$, $x^i_l\in U$, где 


$i\in\{1,2,\dotsc,k,k+1\}$, уже  


определены. Множество 


$U^i_{l+1}=V^i_1\setminus[X<\max\{x^{-1}_lx^{-2}_l,\dotsc, 


x^{-k}_l,x^{k+1}_l\}]$ не 


пусто и открыто в $X$. Выберем $V^i_{l+1}\in P$ и $x^i_{l+1}\in X$ так, 


что  


$[V^{-i}_{l+1}]\subset U^i_{l+1}$ и $x^i_{l+1}\in V^i_{l+1}$. Тогда 


$V^i_{l+1}\in P^*$ и $x^i_{l+1}\in U$ для любого 


$i\in\{1,2,\dotsc,k,k+1\}$. 


Построенные последовательности $\eta^i$ и $\xi^i$ обладают  


свойствами: в)~$[V^i_{l+1}]\subset V$ при $j\in N^+$; 


г)~$x^{i'}_{j'}<y$ для любого $y\in V_{ij}$ при 


$j'<j$ и $i',i\in\{1,2,\dotsc,k,k+1\}$. 


В силу счетной компактности $X$ из~в) следует, что $\eta^i$ сходится  


к $F_i=\cap\eta^i$. Так как $x^i_j\in V^i_j$, то $[\xi^i]\cap 


F_i\ne\varnothing$. Фиксируем $y^i\in F_i\cap[\xi^i]$. 


По~г) $x^i_j<y^i$ при $i',i\in\{1,2,\dotsc,k,k+1\}$, $j'\in N^+$, но 


$y^i\ne y^j$ при  


$i,j\in\{1,2,\dotsc,k,k+1\}$, $i\ne j$, т.к. $y^i\in V^i_1$ и 


$V^i_1\cap V^j_1=\varnothing$ при 


$i,j\in\{1,2,\dotsc,k,k+1\}$, $i\ne j$. Поэтому соотношение  


$[\xi^1\cup\xi^2\cup\xi^k,\,\cup\xi^{k+1}]\supset 


\{y^1,y^2,\dotsc,y^k,y^{k+1}\}$ противоречит  


$k$-согласованности упорядочения $<$ на $X$. 


 


Заметим, что если $x$ --- изолированная точка в $X$, то множество  


$X_{<x}$ замкнуто. 


 


\begin{lem} 


Если $U\in P$ и $|U|>1$, то существуют такие $U'\in P$ и 


$U''\in P$, что $U'\cup U''\subset U$ и $x'<x''$ для любых $x'\in U'$, 


$x''\in U''$. 


\end{lem} 


 


{\bf Доказательство.} Применяя лемму 1 к $U$, получим $V_1\in P$, 


$V_1\subset U$ и $x_1\in U$ такие, что $y<x_1$ для любого $y\in V_1$. 


Если $|V_1|=1$, то  


$U'=V_1$ и $U''=U\setminus X_{<=y}$, где $\{y\}=V_1$, искомые. 


 


Пусть $|V_1|>1$. На следующем шаге применим лемму 1 к $V_1$.  


Получим $V_2\in P$ и $x_2\in V_1$ такие, что $y<x_2$ для любого $y\in 


V_2$. 


Применяя описанную процедуру $k{+}1$ раз, получим $V_{k+1}\in P$ и  


$x_{k+1}\in V_k$ такие, что $y<x_{k+1}$ для любого $y\in V_{k+1}$. При 


$U=V_{k+1}$ имеем  


$x_{k+1}<x_k<\dotsb<x_2<x_1$. Поэтому 


$U''=U\setminus[X_{x_{k+1}}]=\varnothing$. Множества 


$U'$, $U''$ искомые. 


\medskip 


 


{\bf Доказательство теоремы 6.} Если существует $U\in P$, для  


которого $|U|=1$, то заключение теоремы очевидно. Пусть $|U|>1$ 


для любого $U\in P$. Так как $X$ регулярно, то леммы 1 и 2 позволяют  


построить последовательности $\eta=\{U_i:i\in N^+\}$ и $\xi=\{V_i:i\in 


N^+\}$ в $P$ 


такие, что при всех $i\in N^+$:


1)~если $x\in U_i$ и $y\in V_i$, то $x<y$; 


2)~$U_{i+1}\cup[V_{i+1}]\subset V_i$. Положим  


$F=\cap\xi$ и $G=\cup\eta$. Из~2) и счетной компактности $X$ следует, 


что $\xi$ 


сходится к $F$. Но $U_{i+1}\subset V_i$, значит, и $\eta$ сходится к 


$F$. По~1) $x<y$ при 


$x\in G$ и $y\in F$, $W=\cup\{X\setminus[X_{<x}]:x\in F\}$ открыто в 


$X$ и $W\cap G=\varnothing$. 


Так как $\eta$ сходится к $F$, то $F\not\subset W$. Значит, существует 


$x^*=\min F$. 


Пусть $A=[G]\cap F$, тогда $1\le|A|\le K$. Выберем произвольную точку  


$x'\in A$. Далее построим последовательности $\eta'=\{U'_i:i\in N^+\}$ 


и $\xi'=\{V'_i:i\in N^+\}$ в $P$, $V'_i=V_1\setminus(A\setminus\{x'\})$


и $U'_1=U_1$. 


 


Пусть при всех $i<1$ выбраны множества 


$V'_i,U'_i\in P$, которые  


кроме~1), 2) удовлетворяют условию 3)~$x'\in V'_i$. В силу регулярности  


пространства $X$ существует $V\in P$ такое, что $x'\in 


V\subset[V]\subset V'_{l-1}$. Так 


как $x\notin[U_i]$ для любого $i\in N^+$, то существует такое $n\in 


N^+$, что  


$U_n\cap V\ne\varnothing$ и можно считать, что $n\ge 1$. 


 


Множества $U'_l=U_n\cap V$, $V'_l=V_n\cap V\in P$ удовлетворяют  


условиям~1), 2), 3). Итак, последовательности $\xi'$ и $\eta'$ 


построены.  


Положим $F'=\cap\xi'$ и $G'=\cup\eta'$. Как и выше, последовательности  


$\xi'$ и $\eta'$ сходятся к $F$.\; 


$W'\supset F'\setminus\{x'\}$ открыто в $X$ и $W'\cap 


G'\ne\varnothing$. Существует 


$\min F'$ и по условию~3) $\min F'=x'$. В силу выбора множества $V'_i$ 


имеем  


$W'\supset F'\setminus\{x'\}$. Пусть $O_{x'}\in x'$, $O_{x'}\in P$. 


Тогда $O_{x'}\cup W'\supset F'$ и 


поэтому существует такое $i'\in N^+$, что $V_{i'}\subset O_{x'}\cup W'$. 


Значит, по~2)  


и $U'_i\subset O_{x'}\cup W'$ при всех $i>i'$. Но $U'_i\cap 


W'_i=\varnothing$, значит, $U'_i\subset O_{x'}$ 


для любого $i>i'$, т.е. последовательность $\eta'$ сходится к $x'$, 


что и  


требовалось доказать. 


 


В работе [2] введено обозначение $\delta(x,X)\le\aleph_0$, если 


существует  


такое счетное центрированное семейство $\xi$ открытых в $X$ множеств,  


что $x\in\cap\{[U]:\in\xi\}$ и для любой окрестности $O_x$ точки $x$ в 


$X$ 


существует такое $V\in\xi$, что $V\subset O_x$. Из доказательства 


теоремы 6  


получается  


 


\begin{thm} 


Если $X$ принадлежит классу $\Im(k)$, состоящему из 


$k$-растянутых, счетно компактных, регулярных $T_1$-пространств, то  


$[\{x\in X:\delta(x,X)\le\aleph_0\}]=X$. 


\end{thm} 


 


\begin{cor} 


Если $X\in\Im(k)$, то существуют такая  


последовательность $S=\{F_n:n\in N^+\}$ непустых замкнутых в $X$ 


множеств $F_n$ и такая точка $x\in X$, что последовательность $S$ 


сходится  


к точке $x$ и $\aleph(F_n,X)\le\aleph$ для любого $n\in N^+$. 


\end{cor} 


 


\begin{cor} Если $X\in\Im(k)$ и $|X|\ge\aleph$, то в $X$ имеется


нетривиальная сходящаяся последовательность. 


\end{cor} 


 


{\bf Доказательство.} Если $X$ не разрежено, то доказательство  


сводится к случаю, когда в $X$ нет изолированных точек, а тогда 


заключение следует из теоремы 6. Если же $X$ разрежено, то см.  


приложение 3 из работы [2]. 


 


\begin{cor} 


Если $X\in\Im(k)$, то множество точек счетного $\pi$ 


характера в $X$ всюду плотно в $X$. 


\end{cor} 


 


Бикомпакт $X$ называют {\it диадическим}, если он является 


непрерывным образом канторова куба $D^\tau$ при некотором 


$\tau\ge\aleph_0$. 


 


Вполне регулярное пространство называют {\it диадическим}, если 


оно имеет диадическое бикомпактное расширение. 


 


\begin{cor} 


Если $X\in\Im(k)$ и $X$ диадично, то $\omega(X)\le\aleph_0$. 


\end{cor} 


 


Базу топологического пространства называют {\it слабо 


н\"етеровской}, если каждый ее элемент содержится в качестве 


подмножества в не более чем счетном множестве других 


элементов из базы.  


 


Пространство, имеющее слабо н\"етеровскую базу, называют 


{\it слабо н\"етеровским}. 


 


\begin{cor} 


Слабо н\"етеровский $k$-растянутый бикомпакт  


метризуем, следовательно, $\alpha$-растянут.  


\end{cor} 


 


\begin{cor} 


Если бикомпакт $D^\tau$\; $k$-растянут, то $\tau\le\aleph_0$. 


\end{cor} 


 


\begin{cor} 


Каждый $k$-растянутый диадический бикомпакт $X$ 


метризуем, следовательно, $\alpha$-растянут. 


\end{cor} 


 


Следствие 8 справедливо, т.к. иначе $X$ 


содержит топологическую копию $D^\aleph_1$. 


 


\begin{thebibliography}{9} 


 


\bibitem{1} 


Архангельский А.В. О топологиях, допускающих слабую  


связь с упорядочениями // ДАН СССР. -- 1978. -- Т.\,238. -- \No\,4. 


-- С.\,773--776. 


\bibitem{2} 


Архангельский А.В. Об $\alpha$-растянутых пространствах // ДАН  


СССР. -- 1978. -- Т.\,239. -- \No\,3. -- С.\,336--340. 


\bibitem{3} 


Тюгашов Г.С. О $P$-растянутых пространствах // Изв. вузов. Математика. 


-- 1990. -- \No\,10. -- С.\,74--75. 


 


\end{thebibliography} 


 


\vskip 2cm 


 


\post{\it Самарский архитектурно-}{\it Поступила} 


\post{\it строительный институт}{16.09.1994} 


 


\label{end} 


\end{document} 


